\documentclass[11pt,letterpaper]{article}
\usepackage[T1]{fontenc}
\usepackage{lmodern}
\usepackage[margin=0.88in]{geometry}
\usepackage{amsmath,amssymb,amsthm,mathtools}
\usepackage{booktabs,array,longtable}
\usepackage{xcolor,microtype}
\usepackage[hyphens]{url}
\usepackage[colorlinks=true,linkcolor=blue!45!black,citecolor=blue!45!black,urlcolor=blue!45!black]{hyperref}
\hypersetup{pdftitle={Smaller Diophantine Forms for Con(ZFC)},
pdfauthor={Prepared for Vladimir Reshetnikov},
pdfsubject={Pell-root translations, joint positivity, and coefficient-bounded proof certificates},
pdfkeywords={Diophantine equations, ZFC consistency, Pell equations, proof certificates}}
\usepackage{fancyhdr}
\pagestyle{fancy}
\fancyhf{}
\fancyhead[L]{\small Diophantine forms for $\operatorname{Con}(\mathrm{ZFC})$}
\fancyhead[R]{\small Pell translations and coefficient budgets}
\fancyfoot[C]{\thepage}
\setlength{\headheight}{14pt}
\setlength{\emergencystretch}{2em}
\newcommand{\Z}{\mathbb Z}
\newcommand{\N}{\mathbb N}
\newcommand{\Q}{\mathbb Q}
\newcommand{\Pos}{\mathbb Z_{>0}}
\newcommand{\Con}{\operatorname{Con}}
\newcommand{\ZFC}{\mathrm{ZFC}}
\newtheorem{theorem}{Theorem}[section]
\newtheorem{lemma}[theorem]{Lemma}
\newtheorem{proposition}[theorem]{Proposition}
\newtheorem{prop}[theorem]{Proposition}
\newtheorem{corollary}[theorem]{Corollary}
\theoremstyle{definition}
\newtheorem{definition}[theorem]{Definition}
\theoremstyle{remark}
\newtheorem{remark}[theorem]{Remark}

\title{Smaller Diophantine Forms for $\Con(\ZFC)$\\[5pt]
\large Two Pell-root translations, joint positivity, and\\coefficient-bounded proof certificates}
\author{Research note prepared for Vladimir Reshetnikov\\
\small AI-assisted mathematical development with independent algebraic review}
\date{October 23, 2026}

\begin{document}
\maketitle
\begin{abstract}
We construct explicit fixed-coefficient Diophantine equations whose lack of
solutions is equivalent, by ordinary mathematics formalizable in ZFC, to
syntactic $\Con(\ZFC)$. The principal integer constructions have 26 unrestricted
variables at degree 8 and 24 at degree 10; a further restriction of the
dyadic coordinate gives 23 at degree 16. The 24-variable equation is a sum
of twelve squared residuals. Its reduction uses two integral translations
of Pell roots before eliminating coordinates, together with an exact
four-witness representation of two strict inequalities. Two positive-domain
constructions have 11 variables, degree 277 and 111 arithmetic operations,
and 12 variables, degree 78 and 103 operations.

Separately, we prove a new coefficient-budget encoding of finite
single-conclusion proof certificates. Products of linear factors recover
exact proof-label balances when evaluated at a base exceeding an internally
accumulated coefficient bound. This removes the earlier label-difference
common-multiple requirement at the source level. We supply its local
polynomial and tests, but do not claim a completed fixed-dimensional
consistency equation from that architecture. All finished equations retain
the cited native-source calibration. No global minimum, unconditional
world-first claim, or new complete proof-assistant certificate is asserted.
\end{abstract}

\tableofcontents

\section{Results, comparison, and the exact logical claim}
\label{sec:scope}

The input was the current ProveIt collection
\path{Computability/DiophantineConsistency/con-zfc-diophantine-forms}.
The reviewed checkout is pinned to commit
\begin{center}
\texttt{2605669a532342f153b13eb6b4a0af4b54f5fa37}.
\end{center}
Its latest package reconciliation includes sources 68--75, the tenth wave. The review does
not compare with an obsolete earlier wave. The repository already has
nine-positive-variable and ten-integer equations of much higher degree,
as well as a 38-integer quartic, a 32-integer sextic and a 27-integer
octic\cite{frontier,source68,source69,source70,source71}.

There is no agreed total ordering of simplicity. Variable count, domain,
degree, circuit size, coefficient height and proof complexity measure
different things. The present results improve several precise parts of
that comparison; they do not supersede every earlier candidate.

For broader context, the current general representation literature includes
Sun's ten-integer result and the parallel formalization work of Bayer,
David, Ha\ss ler, Matiyasevich and Schleicher\cite{sun2026,bayer2026}.
Those universal bounds do not by themselves give a smaller low-degree
specialization for this fixed consistency predicate. The positive
representation and Pell methods used below also lie in the framework
of Sun's earlier work\cite{sun2021}.

\begin{center}
\small
\begin{tabular}{@{}lrrrl@{}}
\toprule
Construction & Variables & Degree & Operations & Comparison\\
\midrule
Integer octic & 26 & 8 & -- & 27-variable degree-8 predecessor\\
Integer decic & 24 & 10 & -- & new variable/degree tradeoff\\
Integer degree-16 companion & 23 & 16 & -- & 26-variable degree-16 predecessor\\
Positive equation $P_{11}$ & 11 & 277 & 111 & predecessor: 325 / 113\\
Positive equation $P_{12}$ & 12 & 78 & 103 & predecessor: 82 / 102\\
\bottomrule
\end{tabular}
\end{center}

An operation is a live binary addition, subtraction or multiplication;
squares and multiplication by constants are charged. Fixed integer leaves
and reuse of an already computed register are free in this convention.
The integer systems are delivered as explicit residual lists and sparse
expansions; their evaluator counts have not been optimized or compared.
The last positive result trades one operation for four degrees.

\subsection{Consistency convention and dependencies}

Fix the recursively specified syntactic ZFC proof predicate used by the
repository. Write $\operatorname{Ref}_{\ZFC}(p)$ for the statement that $p$
is a finite refutation, and
\[
 \Con(\ZFC)\equiv\neg\exists p\in\N\;\operatorname{Ref}_{\ZFC}(p).
\]
The inherited native-source theorems identify the existence of their
Diophantine witnesses with the existence of such a refutation. Those
theorems include the Wang-tile/native compiler, Pell index reconstruction,
and the fixed machine/calculus calibration. We use these as named
mathematical inputs and preserve their coefficient conventions.

For each finished equation $P$ with its stated domain $D$, the result is
\begin{equation}
 \ZFC\vdash
 \bigl[\neg\exists\boldsymbol{x}\in D^n\;P(\boldsymbol{x})=0\bigr]
 \;\longleftrightarrow\;\Con(\ZFC).
 \label{eq:target}
\end{equation}
The new proofs establish both existence implications without assuming
consistency. They then compose with the cited source calibration and
negate the equivalence. Thus the argument is neither a bounded search
for contradictions nor a proof that ZFC is consistent.

Equation~\eqref{eq:target} is an ordinary mathematical provability claim.
It is stronger than an assertion merely about a few computed witnesses.
It is also distinct from delivering a checked term in the repository's
internal first-order ZFC calculus. The latter development remains
incomplete for these native-source endpoints\cite{objectzfc}. This archive
adds explicit equations, reductions and exact algebra checks; it does not
label finite Python checks as a formal proof of the unbounded semantic
calibration.

\subsection{What is new in this note}

The nearest-square positivity lemma and determinant guard are inherited
from different tenth-wave packages. Their composition gives the 26-variable
octic. The more consequential algebraic step is to translate \emph{both}
Pell roots before removing the $c$ and $V$ coordinates: this yields the
24-variable degree-10 equation while keeping all eliminations integral.
The degree-16 companion uses a separately justified divisibility restriction
on the canonical dyadic coordinate.

The positive-domain changes likewise combine results that were previously
separate: the adjacent-coefficient descent makes a denominator-protection
factor redundant, and the adaptive Pell carrier lowers a second degree.
The coefficient-budget proof source is a further mechanism, developed
here from the previous commutative-flow proposal. Its theorem is complete;
its compilation into a competitive whole consistency polynomial is not.

\section{Joint guards and two Pell-root translations}
\label{sec:integer-new}

This section gives three explicit fixed-coefficient equations over unrestricted
integers.  The principal constructions have 26 variables and degree eight,
and 24 variables and degree ten.  A further restriction to canonical dyadic
witnesses gives 23 variables and degree sixteen.  They improve the inspected
repository's 27-variable octic and 26-variable degree-sixteen construction
\cite{source68,source71}.  No global minimum or priority claim outside the
inspected repository is made.

The starting point is the literal 38-variable quartic of source 71.  Its
fixed tuple and determinant guard are kept intact.  Two operations are
combined here: the simultaneous-positivity lemma of source 68 is applied
to the two remaining guards, and two integral translations of Pell-root
coordinates permit further substitutions without the usual degree increase.
The first operation is a synthesis of earlier results.  The paired
translations and the resulting 24-variable system are the additional
algebraic reduction.

\subsection{One guard for two strict inequalities}

Put
\[
 \operatorname{Pos}(p_1,p_2,p_3)
   =1+p_1^2+p_2^2+p_3(p_3+1).
\]
This takes only positive integer values and represents every positive
integer.  Indeed, apply the three-square theorem to $4n-3$ and separate its
one odd coordinate from its two even coordinates.  This gives
$4n-3=(2p_3+1)^2+(2p_1)^2+(2p_2)^2$.

\begin{lemma}[Nearest-square simultaneous positivity; source 68]
\label{lem:int-pair}
For any $x,y\in\Z$,
\[
 x>0\ \land\ y>0
 \quad\Longleftrightarrow\quad
 \exists t,p_1,p_2,p_3\in\Z:
 xy-(t^2-x-y+1)^2=\operatorname{Pos}(p_1,p_2,p_3).
\]
\end{lemma}
\begin{proof}
A solution makes $xy>0$, so $x,y$ have the same nonzero sign.  If
$x=-a$, $y=-b$, where $a,b\ge1$, then
$(t^2+a+b+1)^2>ab=xy$, a contradiction.

Conversely let $x,y\ge1$, put $N=x+y-1$, and take
$m=\lfloor\sqrt N\rfloor$.  Choose $t=m$ if $N-m^2\le m$, and
$t=m+1$ otherwise.  If $N$ is a square the error is zero.  Otherwise
$|t^2-N|\le m$ and $m^2<N$.  Hence
\[
 (t^2-N)^2<N\le xy,
 \qquad xy-N=(x-1)(y-1)\ge0.
\]
The remaining positive integer has a $\operatorname{Pos}$ representation.
The floor and square root are used only to choose a witness in this proof;
they are not operations in the polynomial.
\end{proof}

Two separate strict inequalities previously cost six private integer
coordinates.  Lemma~\ref{lem:int-pair} uses four.  When $x$ and $y$ are
quadratic polynomial expressions, its residual has degree at most four.
It is essential to keep both inputs quadratic while forming the octic:
eliminating $\zeta$ too early changes $\eta\zeta$ into a cubic expression.

\subsection{Fixed coefficients}

The six fixed integers are
\[
 \beta=2^{1064921}-1,\quad K_0,\quad M_C,\quad M_F,
 \qquad \sigma_F=2^{1113},\quad \sigma_C=2^{666104}.
\]
Here $(\beta,K_0,M_C,M_F)$ is the complete
\texttt{mixed-parity-unrounded-1064921} tuple of source 34, used without
mixing ports from other tuples.  In particular $M_F$ is already shifted
by $\beta$.  The archive supplies every actual integer as compressed
minimal unsigned little-endian bytes, its SHA-256 digest, and the finite
recipe rebuilding it from the literal K215 transition table and ordered
943-tile list.  Thus these are fixed numerals, not parameters to be chosen
after searching for a solution.  The coefficient proof and calibration
remain those of sources 34, 64 and 71 \cite{source64,source71}.

\begin{center}
\begin{tabular}{lrr}
\hline
Fixed numeral & Bit length & Minimal raw bytes\\
\hline
$\beta$ & 1,064,921 & 133,116\\
$K_0$ & 1,065,043 & 133,131\\
$M_C$ & 1,064,921 & 133,116\\
$M_F$ & 1,064,922 & 133,116\\
$\sigma_F$ & 1,114 & 140\\
$\sigma_C$ & 666,105 & 83,264\\
\hline
Total & & 615,883\\
\hline
\end{tabular}
\end{center}

\subsection{The octic in 26 integers}

Its independent integer coordinates are
\[
\begin{gathered}
 J,F,\alpha,z,f,h,r,v,w,\tau,\eta,\zeta,y,\gamma,o,j,\\
 E,c,a,V,I,u,t,p_1,p_2,p_3.
\end{gathered}
\]
The code names $v,w,t,p_i$ as \texttt{v\_scale}, \texttt{w\_scale},
\texttt{pair\_root}, and \texttt{guard1},\ldots,\texttt{guard3}.
Every expression in the following display is an abbreviation for a
polynomial in these 26 coordinates:
\begin{align*}
 q&=\beta J+1, & k&=\eta+\zeta, & Y&=a-E,\\
 C&=q-F-\alpha, & X&=2u^2, & H&=4a+3,\\
 \delta&=a^2+H, & T&=E(c-\eta), & S&=X+\gamma H,\\
 R&=\bigl((q-1)(q-F)+\alpha\bigr)(q^2-1)
       +(M_C+qM_F)J.&&&
\end{align*}
The two guard expressions and their combined residual are
\begin{align*}
 G&=(\sigma_F F)(\sigma_C C)
       -(q-\sigma_F F-\sigma_C C)^2,\\
 Z&=\eta\zeta,\\
 B_*&=GZ-(t^2-G-Z+1)^2-\operatorname{Pos}(p_1,p_2,p_3).
\end{align*}
Define the fourteen residuals
\begin{align}
 A_1&=X-4qw, & A_2&=a-E-q^3(2v+1),\nonumber\\
 A_3&=E-XY, & A_4&=c-Yk-\eta,\nonumber\\
 A_5&=\tau^2+(2\tau-k)T-1,&&\nonumber\\
 A_6&=2acS+S^2-Hc^2-1,&&\nonumber\\
 A_7&=k-1-hE-R,&&\nonumber\\
 A_8&=(K_0+4w)C-(q-1)z-F,&&\nonumber\\
 A_9&=V+c-fo, & A_{10}&=V+k-1-cj-hE,\nonumber\\
 A_{11}&=I-(2r+1)c^2, & A_{12}&=I^2-\delta(f^2-1),\nonumber\\
 A_{13}&=I^2(V^2-y^2)+y^2-1, & A_{14}&=B_*.
 \label{eq:int26rows}
\end{align}
The requested single equation is
\begin{equation}
 \boxed{\displaystyle Q_{26}=\sum_{i=1}^{14}A_i^2=0.}
 \label{eq:int26}
\end{equation}
All quantifiers range over $\Z$.  The abbreviations introduce no additional
quantified coordinates and contain only fixed nonnegative powers.

\begin{proposition}
The equation $Q_{26}=0$ has an integer solution exactly when the literal
source-71 quartic has an integer solution.  Its exact total degree is eight.
\end{proposition}
\begin{proof}
The source quartic has two rows $G=\operatorname{Pos}$ and
$Z=\operatorname{Pos}$, each with its own private triple.  Replace these
six coordinates and two rows by $t,p_1,p_2,p_3$ and $B_*=0$.
Lemma~\ref{lem:int-pair} proves equivalence on the projection to all the
other coordinates.  This gives a 36-variable system of 24 residuals.

Translate the first root by
\[
 \tau_{\rm old}=\tau+E(c-\eta).
\]
Then eliminate the following ten monically defined coordinates:
\[
\begin{array}{c|c@{\qquad}c|c}
 X&2u^2 & C_2&c^2\\
 F_2&f^2 & B_{\rm aux}&I^2\\
 g&y^2-V^2 & \Delta&a^2+4a+3\\
 D&2u^2+ac+\gamma(4a+3)&P&k+E(c-\eta)\\
 L&q^2-1&N&q^3
\end{array}
\]
These are integer-polynomial restorations with integer-polynomial inverses
on the zero sets, given by forgetting the defined coordinates and undoing
the root translation.  The ten definition rows become identities.
For the first and main norms the exact identities are
\begin{align*}
 (\tau+T)^2-(k+T)^2+(k+T)k-1
   &=\tau^2+(2\tau-k)T-1,\\
 (ac+S)^2-(a^2+H)c^2-1
   &=2acS+S^2-Hc^2-1.
\end{align*}
The surviving rows are exactly \eqref{eq:int26rows}, up to multiplication
of individual residuals by $-1$, which leaves their squares unchanged.
The literal source-row restorations are also checked by exact sparse
polynomial arithmetic in the archive.

Every residual has degree at most four.  The only residual containing
$t$ is $B_*$, whose $t^4$ coefficient is $-1$.  Therefore
$[t^8]Q_{26}=1$, independently of the fixed coefficients.  This proves
exact degree eight on the actual coefficient slice.
\end{proof}

\subsection{The degree-ten equation in 24 integers}

Retain the same definitions, but now set
\begin{equation}
 c=Yk+\eta,\qquad V=fo-c,
 \qquad y_{\rm old}=V+\nu.
 \label{eq:int24changes}
\end{equation}
Here $c,V$ are polynomial abbreviations.  Replace the coordinate $y$ by
$\nu$ and remove the coordinates $c,V$.  The 24 independent coordinates are
\[
\begin{gathered}
 J,F,\alpha,z,f,h,r,v,w,\tau,\eta,\zeta,\nu,\gamma,o,j,\\
 E,a,I,u,t,p_1,p_2,p_3.
\end{gathered}
\]
Remove $A_4,A_9$, which now vanish identically.  Replace $A_{13}$ by
\begin{equation}
 A'_{13}=(V+\nu)^2-I^2\nu(2V+\nu)-1.
 \label{eq:int24aux}
\end{equation}
In every other residual substitute the two definitions in
\eqref{eq:int24changes}.  The complete single equation is
\begin{equation}
 \boxed{\begin{aligned}
 Q_{24}={}&A_1^2+A_2^2+A_3^2+A_5^2+A_6^2+A_7^2\\
         &+A_8^2+A_{10}^2+A_{11}^2+A_{12}^2+(A'_{13})^2+B_*^2=0.
 \end{aligned}}
 \label{eq:int24}
\end{equation}

\begin{theorem}[Two root translations and two coordinate eliminations]
The integer zero sets of $Q_{26}$ and $Q_{24}$ are related by inverse
integer-polynomial maps.  The polynomial $Q_{24}$ has exactly 24 live
integer variables and exact total degree ten.
\end{theorem}
\begin{proof}
The two definition rows $A_4,A_9$ are monic in $c,V$.  Their right-hand
sides in \eqref{eq:int24changes} use only retained coordinates and have
degree two.  The auxiliary-root translation has the integral inverse
$\nu=y_{\rm old}-V$.  The exact identity
\[
 I^2\bigl(V^2-(V+\nu)^2\bigr)+(V+\nu)^2-1
  =(V+\nu)^2-I^2\nu(2V+\nu)-1
\]
proves restoration of the auxiliary norm for arbitrary signed integers.
There are no sign assumptions or divisibility conditions in these maps.

The degree savings require both root translations.  After eliminating
$c$, the expression $T=E(c-\eta)$ has degree three.  Centering the first
root makes its norm $\tau^2+(2\tau-k)T-1$ of degree four instead of six.
After eliminating $V$, the auxiliary-root translation in
\eqref{eq:int24aux} makes its norm degree five instead of six.  In the
main norm the cancellation of $a^2c^2$ already displayed leaves
$2acS+S^2-Hc^2-1$, also of degree five.  Finally
$I-(2r+1)c^2$ has degree five.  All other residuals have degree at most
four.  In the order displayed in \eqref{eq:int24}, their degrees are
\[
 2,4,3,4,5,4,2,3,5,4,5,4.
\]
Thus $\deg Q_{24}\le10$.

Only $A_{11}$ contains $r$.  The top part of
$c=(a-E)(\eta+\zeta)+\eta$ gives
\[
 [r^2a^4\eta^4]Q_{24}=4.
\]
This coefficient is independent of all six fixed numerals, so the degree
is exactly ten after actual coefficient specialization.  Exact symbolic
checks also verify that every listed coordinate occurs.
\end{proof}

\subsection{An optional 23-variable equation of degree sixteen}

Starting from $Q_{24}$, substitute simultaneously
\begin{equation}
 u_{\rm old}=2q\,u,\qquad w=2q\,u^2.
 \label{eq:int23dyadic}
\end{equation}
Then $X=8q^2u^2$ and $A_1=X-4qw$ is identically zero.  Remove $w$ from
the coordinate list and remove this zero residual.  The other eleven
squared residuals define $Q_{23,16}$, a polynomial in 23 unrestricted
integer variables, with the same six fixed coefficients.

\begin{proposition}
The absence of integer zeros of $Q_{23,16}$ has the same consistency
equivalence as $Q_{24}$.  Its exact degree is sixteen.
\end{proposition}
\begin{proof}
Any new zero restores a $Q_{24}$ zero by the integral substitutions
\eqref{eq:int23dyadic}.  Completeness uses fresh canonical native witnesses,
rather than asserting that every old zero satisfies the restriction.
Source 64's canonical witnesses have
\[
 q=2^n,\quad n\ge4,\qquad R\text{ odd},\quad R>3q,
 \qquad u_{\rm old}=2^{(R-1)/2}.
\]
Since $3\cdot2^n>2n+3$ for $n\ge4$, the exponent
$(R-1)/2-n-1$ is a nonnegative integer.  Thus $2q$ divides
$u_{\rm old}$, and the new $u$ is integral.  The old equation
$2u_{\rm old}^2=4qw$ then gives exactly $w=2qu^2$.
The source-71 determinant construction retains those canonical coordinates,
and the new joint guard and root translations do not change this fact.

With $X=8q^2u^2$, the main-norm residual has degree eight.  Every other
residual has degree at most five.  Its unique top contribution
$X^2=64q^4u^4$ gives
\[
 [J^8u^8]Q_{23,16}=4096\beta^8>0.
\]
Therefore the exact degree is sixteen.  This step is an existence
equivalence obtained through canonical witnesses, not a bijection of all
old zero sets.
\end{proof}

\subsection{Consistency theorem and verification boundary}

Composing the proved transformations with the named source-71 correctness
theorem yields, by ordinary arguments formalizable in ZFC, for each of
$Q_{26},Q_{24},Q_{23,16}$,
\[
 \mathrm{ZFC}\vdash
 \bigl(\neg\exists\mathbf{x}\in\Z^m\;Q(\mathbf{x})=0\bigr)
 \;\longleftrightarrow\;\operatorname{Con}(\mathrm{ZFC}).
\]
The consistency predicate is the specified syntactic ZFC refutation
predicate of the inherited calibration.  The implication does not assert
either consistency or inconsistency of ZFC.  No variable coefficient,
variable exponent, hidden inequality, or quantified abbreviation is used.

The release checker reconstructs all six actual fixed integers, checks
their complete raw hashes, verifies every restored residual exactly,
expands each sum of squares over the coefficient ports, proves the
reported degree by the displayed nonzero coefficient, and checks variable
liveness.  Its finite signed and nearest-square samples are only regression
diagnostics; Lemma~\ref{lem:int-pair} is the universal proof.  The inherited
unbounded native compiler, Pell classification and rank arguments, and
literal machine/refutation calibration are not replaced by those finite
checks.  This package adds no end-to-end proof-assistant or internally
encoded ZFC certificate.

\section{Two positive-domain improvements by combining the wave-10 mechanisms}
\label{sec:positive-improvements}

The integer-domain construction in the preceding sections is not the only
useful endpoint.  The same review yields two improvements to the positive
frontier.  They are separate, fully specified polynomials with the authentic
compact44 coefficient tuple.  Their contribution is a precise composition
of the mechanisms in sources~69 and~70, rather than a claim that the generic
norm or descent methods are new.

\begin{center}
\begin{tabular}{lrrr}
\hline
Construction & Positive variables & Exact degree & Operations\\
\hline
Source 70 & 11 & 325 & 113\\
Guard-free strong ordinate & 11 & 277 & 111\\
Source 69 & 12 & 82 & 102\\
Adaptive main carrier & 12 & 78 & 103\\
\hline
\end{tabular}
\end{center}

The operation convention counts every live binary addition, subtraction,
and multiplication on witnesses, including squaring and scalar
multiplication. Fixed integer leaves and register reuse are free. These are
arithmetic-circuit counts, not bit-complexity estimates.

\subsection{An adjacent-coefficient descent}

We use the following arithmetic lemma from source~69\cite{source69}.
Its role here is to remove a hypothesis that source~70 enforced by an
additional factor in its strong Pell ordinate.

\begin{lemma}[Source 69: rational ordinate descent]
\label{lem:adjacent-descent}
Let $Q>1$ and $V$ be integers. If $y\in\mathbb Q$ and
\[
 QV^2-(Q-1)y^2=1,
\]
then $y\in\Z$.
\end{lemma}
\begin{proof}
Write $y=a/m$ in lowest terms, with $m>0$. The equation implies
$m^2\mid Q-1$, and excludes $V=0$ and $y=0$. Change signs to make both
positive. If $V=1$ then $y=1$. Otherwise define
\[
 V'=(2Q-1)V-2(Q-1)y,\qquad
 y'=(2Q-1)y-2QV.
\]
The first is integral because $m\mid Q-1$. Direct expansion preserves the
norm. The reduced denominator of $y'$ is exactly $m$, since
$2Q-1\equiv1\pmod m$. Moreover $y>V$ and $QV>(Q-1)y$, and hence
$-V<V'<V$. The norm excludes $V'=0$. Strong induction on the positive
integer $V$, after taking absolute values, reaches the case $V=1$
without changing the denominator. Thus $m=1$.
\end{proof}

The exact polynomial norm identity in this proof is checked by the supplied
twelve-variable verifier. No finite enumeration is being used in place of
the induction.

\subsection{Eleven positive unknowns in degree 277}

Use the eleven inputs and every register of source~70's eleven-positive
construction\cite{source70}, with exactly one change:
\[
 \boxed{I=ic^2\quad\hbox{in place of}\quad I=iq(q-1)c^2.}
\]
The inputs are
\[
 (j,F,\alpha,i,T,X,Y,n,\eta,\zeta,\gamma)\in\Pos^{11}.
\]
The fixed numerals are $\beta,K_0,M_C,M_F$, with
$\beta=2^{590069}-1$ and the already-shifted $M_F$.
All abbreviations below are polynomial registers, not extra unknowns.
For completeness the definition is
\begin{align*}
 J&=3j+1,&q&=\beta J+1,&C&=q-F-\alpha,\\
 E&=XY,&k&=\eta+\zeta,&c&=kY+\eta,\\
 a&=E+Y,&H&=4a+3,&\Delta&=a^2+H,\\
 \ell&=EkY,&I&=ic^2,&M&=q^3(q-1),\\
 R&=((q-1)(q-F)+\alpha)(q^2-1)+(M_C+qM_F)J,\\
 R_+&=q^3(q+F+\alpha),&
 B_t&=(qK_0+X)C-qF,&D_\tau&=\ell(\ell+k)+1,\\
 D_f&=1+\Delta I^2,&Q&=\Delta(D_f-1),&\delta&=Q-1,\\
 V&=D_f(cT-R)-c,&L&=3D_f(cT+R_+),&W&=1+(q-1)L,\\
 Z&=M(\ell+k+n+L)+(q-1)Y-q^3WB_t,\\
 S&=4qY,&U&=qS,&v&=k-R-1,\\
 D_*&=U(X+ac)+H(U\gamma+SX-v),\\
 A_1&=D_*^2-U^2(1+\Delta c^2).
\end{align*}
Now put
\begin{align*}
 A&=M^2D_\tau,&B_f&=M^2D_f,&C_y&=M^2(QV^2-1),\\
 N_c&=\delta^2\bigl[(Z^2-A-B_f)^2-4AB_f\bigr]
       -2\delta C_y(Z^2+A+B_f)+C_y^2,\\
 N_8&=N_c^2-64\delta^3AB_fC_yZ^2.
\end{align*}
The required equation is
\begin{equation}
 \boxed{P_{11}:=(1+N_8)(1+A_1)-1=0.}
 \label{eq:positive277}
\end{equation}
The exact division-free definition is also supplied as
\texttt{positive/eleven/p11\_positive\_degree277.json}.

\begin{prop}
ZFC proves that \eqref{eq:positive277} has no zero in $\Pos^{11}$ if and
only if $\operatorname{Con}(\mathrm{ZFC})$. Its exact total degree is 277
and its emitted evaluator has 111 live binary operations.
\end{prop}
\begin{proof}
We give the complete change to the upstream proof; its normalized Pell
and native compiler interfaces remain the ones in sources~61 and~70
\cite{source61,source70}. The source~70 inequalities use only $I\ge c^2$,
which still holds. In particular, writing
\[
 D_y=\frac{QV^2-1}{Q-1},
\]
the three positive rational radicands $D_\tau,D_f,D_y$ are pairwise
distinct, and no signed sum of their square roots is zero. The octic norm
therefore recovers rational signed roots $\tau_*,f_*,y_*$ with squares
$D_\tau,D_f,D_y$. The first two are integers because their squares are
integers. Lemma~\ref{lem:adjacent-descent} makes the third an integer.
Here $Q,V$ are integers and $Q>1$ before any equation is imposed.

This replaces the old argument that the denominator of $y_*$ divides $M$
and is coprime to $Q-1$. That argument was the only reason for the extra
$q(q-1)$ in $I$. The unchanged recovery formulas are
\[
 f_{\rm old}=|f_*|,\quad y_{\rm old}=|y_*|,\quad
 u=y_{\rm old}-V>0,\quad r=n+L-f_*-y_*>0.
\]
The signed roots must be used in the formula for $r$. The source proof's
bound $L>|f_*|+|y_*|$ makes it positive. The strong ordinate $I=ic^2$
has exactly the divisibility needed by the normalized index theorem.
Consequently the original four-residual positive source interface is
recovered, and its calibration yields a finite syntactic ZFC refutation.

Conversely any positive zero of source~70 maps directly to a zero here:
replace its input $i$ by $q(q-1)i$ and leave every other input unchanged.
This preserves $I$ and every downstream register. Hence source~70's
completeness supplies completeness here with no new index restriction.
The unit-product finalizer is unchanged: both $N_8$ and $A_1$ are squares
modulo 4, so neither can equal $-2$.

For exact degree, the changed register degrees are
\[
\begin{array}{c|rrrrrrrrr}
\text{register}&I&D_f&Q&L&W&Z&N_c&N_8&A_1\\\hline
\text{degree}&5&14&18&18&19&24&132&264&13.
\end{array}
\]
Writing $q_0=3\beta j$, $k_0=\eta+\zeta$ and
$B_0=(K_0q_0+X)(q_0-F-\alpha)-q_0F$, the leading forms are
\begin{align*}
 Q_0&=i^2k_0^4X^4Y^8,\\
 Z_0&=-3i^2q_0^7k_0^4X^2Y^6(q_0+F+\alpha)B_0,\\
 (A_1)_{13}&=32q_0^2k_0X^2Y^4
      \bigl(4q_0^2Y\gamma+q_0^3(q_0-F)\bigr).
\end{align*}
The unique top term of $N_8$ is $Q_0^4Z_0^8$. Therefore the coefficient
of
\[
 j^{60}i^{24}\eta^{49}X^{42}Y^{85}\alpha^{16}\gamma
\]
in $P_{11}$ is $2^7 3^8(3\beta)^{60}>0$. The monomial has degree 277.
The exact checker verifies these identities and the literal graph count:
65 multiplications, 29 additions and 17 subtractions.
\end{proof}

\subsection{Twelve positive unknowns in degree 78}

For a second endpoint, keep all source~69 twelve-variable registers
\cite{source69}, except replace its main Pell carrier by
\begin{align*}
 S&=4qY, & U&=qS=4q^2Y,\\
 D_*&=U(X+ac)+H\bigl(U\gamma+SX-(k-R-1)\bigr),\\
 A_1&=D_*^2-U^2(1+\Delta c^2).
\end{align*}
The source's first carrier $A_0$, two-root resultant $N$, and finalizer
are unchanged:
\begin{equation}
 \boxed{P_{12}=(1+A_0)(1+A_1)(1+N)-1=0.}
 \label{eq:positive78}
\end{equation}
Its full literal definition is
\texttt{positive/twelve/equation.json}; the twelve positive inputs are
$j,F,\alpha,t,T,X,Y,r,\eta,\zeta,n,\gamma$.

To make the unchanged part explicit, use $J=6j+1$, $q=\beta J+1$,
$C=q-F-\alpha$, and the same $E,k,c,a,H,\Delta,\ell,R,B_t$ formulas
as above. Put
\begin{align*}
 W_0&=1+(q-1)k,& M&=q^3(q-1),\\
 Z_0&=M(\ell+k+r)+(q-1)Y-q^3W_0B_t,\\
 A_0&=Z_0^2-M^2\bigl(\ell(\ell+k)+1\bigr),\\
 f&=t+1,& \rho&=f^2-1,& d&=\Delta c^4,\\
 Q&=\Delta\rho,& V&=f^2(cT-R)-c,\\
 h(z)&=(1-Q)z^2+2(1-Q)Vz+V^2-1,\\
 A_N&=dh(n)+(1-Q)\rho,& B_N&=2(1-Q)(n+V),\\
 N&=A_N^2-dB_N^2\rho.
\end{align*}
All coefficients are fixed integers and all powers have fixed exponents.

\begin{prop}
ZFC proves that \eqref{eq:positive78} has no zero in $\Pos^{12}$ if and
only if $\operatorname{Con}(\mathrm{ZFC})$. Its exact total degree is 78
and its emitted evaluator has 103 live binary operations.
\end{prop}
\begin{proof}
The unit-product finalizer still works: $q$ is even, $4\mid U$, and all
three residuals are squares modulo 4. First $A_0=0$ gives $q^3\mid Y$
and $C>0$, hence $q\ge16$, $q\equiv2\pmod3$ and
$3q<R<q^4-2$. The new $D_*$ is positive, since
$D_*>Uac-Hk>ak(UY-5)>0$. Thus $A_1=0$ recovers the positive main
Pell root. The first and main roots give $c>2R$ before the strong or
auxiliary equations are used. Consequently $V>1$, and the unchanged
resultant and Lemma~\ref{lem:adjacent-descent} recover positive integral
strong and auxiliary coordinates. The normalized auxiliary-index theorem
then gives $p=R$.

For the altered carrier set $\mu=S/\gcd(S,H)$. Since $q^3\mid Y$,
$q\equiv2\pmod3$ and $H=4Y(X+1)+3$, one has $\gcd(S,H)\le3$.
Thus $\mu\ge4qY/3>q^4>R+1$. Both $S$ and $U$ divide $4XY^2$
without assuming $q\mid X$. Carrier divisibility gives
$\mu\mid k-R-1$, and the first Pell recurrence gives
$k\equiv2n_{\rm first}\pmod{4XY^2}$. With
$1\le n_{\rm first}<p=R<\mu$, these imply
$2n_{\rm first}=R+1$. Therefore $v=k-R-1$ is a positive multiple of
$4XY^2$, hence of $U$. Returning to the carrier gives $q\mid HX$,
and $\gcd(q,H)=1$ gives $q\mid X$.

Writing $w=X/q$, the restored source gap is
$\gamma_{\rm old}=\gamma+w-v/U$. It is integral and positive because
the main root is greater than $(a+1)c$ and $c>X$.
This recovers the source~69 interface. Conversely its canonical witnesses
lift by $\gamma=\gamma_{\rm old}+v/U-w>0$, using
$v/U\ge XY/q^2=wY/q\ge wq^2>w$. No new canonical restriction is
introduced.

The exact residual degrees are $15,13,50$. Their product has degree 78:
its coefficient at
\[
 j^{18}F^2\eta^{11}X^{12}Y^{23}t^8n^2\gamma\alpha
\]
is $1024(6\beta)^{18}>0$. The supplied checker verifies the complete
carrier expansions, resultant identities and this coefficient. It counts
58 multiplications, 26 additions and 19 subtractions. Thus this endpoint
trades one operation for four degrees compared with source~69.
\end{proof}

These two improvements preserve the native compiler and syntactic
consistency calibration. Their exact algebra is independently replayable,
but they are not presented as new end-to-end proof-assistant certificates.

\section{Proof certificates with linear factors and a coefficient budget}
\label{sec:coefficient-budget}

This section develops a different arithmetic source for proof existence.
The proved contribution is a certificate theorem and its literal local
polynomial.  The final elimination of the unbounded trace has not been
materialized or measured.  Thus the local counts below are not a claim of a
new complete Diophantine-variable record.

Source~68 of the repository\cite{source68} already replaces ordered proofs by nonnegative
Horn flow and encodes formula labels with the pairwise coprime numbers
$1+Mj$.  It requires a second bounded universal block ensuring that $M$ is
divisible by all integers up to a label bound.  We remove that requirement.
We use the linear factors $X+j$, respect the actual arity of each rule, and
carry a separate, small arithmetic budget bounding polynomial coefficients.
The budget is chosen before the evaluation point.  No factorial, primality
test, common-multiple test, or variable exponent occurs in the new source.
The later elimination of bounded universal quantification may still use such
graphs; their costs have not disappeared from a final compiler by assertion.

\newtheorem{budgetlemma}{Lemma}[section]
\newtheorem{budgettheorem}[budgetlemma]{Theorem}
\newtheorem{budgetcorollary}[budgetlemma]{Corollary}

\subsection{Exact evaluation detects a polynomial identity}

All formula labels are natural numbers, including zero.  A rule has one
conclusion and zero, one, or two premises.  An occurrence of zero as an
actual premise is a genuine obligation.  In the C222 application, zero is
discharged by the explicit truth axiom; it is not discarded as a free
coordinate in the following construction.

For a finite list of rule instances
$R_i=(P_i\longrightarrow c_i)$, $0\leq i<N$, put
\begin{align}
 B_i&=\prod_{p\in P_i}(1+p)+c_i,
 &T_0&=1,
 &T_{i+1}&=T_i B_i,                                      \label{eq:budget-update}\\
 C(X)&=\prod_{i<N}(X+c_i),
 &Q(X)&=\prod_{i<N}\prod_{p\in P_i}(X+p),
 &M&=T_N+1.                                               \label{eq:budget-polynomials}
\end{align}
Products count repeated premises with their multiplicities.  An empty
product is one.  If two numerical premise slots $p_i,q_i$ are used, pad
unused slots with zero.  Then
\[
 B_i=(1+p_i)(1+q_i)+c_i
\]
for a binary rule; the same expression specializes to $1+p_i+c_i$ for a
unary rule and to $1+c_i$ for an axiom.

\begin{budgetlemma}[Coefficient budget]
Every coefficient of $C(X)-(X+1)Q(X)$ has absolute value at most $T_N$.
Consequently,
\begin{equation}
 C(M)=(M+1)Q(M)
 \quad\Longleftrightarrow\quad C(X)=(X+1)Q(X)\text{ in }\mathbb Z[X].
 \label{eq:budget-evaluation}
\end{equation}
\end{budgetlemma}
\begin{proof}
Each $B_i$ is at least both $1+c_i$ and
$\prod_{p\in P_i}(1+p)$.  Thus
$C(1)\leq T_N$ and $Q(1)\leq T_N$.
Write $Q(X)=\sum_j q_jX^j$, with $q_j\geq0$ and missing coefficients zero.
The coefficient of $X^j$ in $(X+1)Q(X)$ is $q_j+q_{j-1}$.
These are distinct entries in the coefficient list, so their sum is at
most $Q(1)$, including at either endpoint.  Every coefficient of $C$ is
at most $C(1)$.  The difference of two nonnegative numbers each at most
$T_N$ has absolute value at most $T_N$, not $2T_N$.

It remains to use a simple injectivity argument.  If a nonzero integer
polynomial $H(X)=\sum h_jX^j$ has $|h_j|<M$, let $k$ be the least index
with $h_k\ne0$.  Dividing $H(M)=0$ by $M^k$ and reducing modulo $M$ gives
$M\mid h_k$, contradicting $0<|h_k|<M$.  Apply this to
$H=C-(X+1)Q$ and $M=T_N+1$.  The reverse implication is immediate.
\end{proof}

The factor two sometimes used in such a coefficient bound is unnecessary.
Also, the product
$\prod_i(1+c_i)(1+p_i)(1+q_i)$ is an unnecessarily large budget.
The factor $B_i$ in~\eqref{eq:budget-update} gives the required two
majorizations separately and lowers the budget update from degree four to
degree three in its displayed coordinates.

\begin{budgetlemma}[Linear-factor uniqueness]
The polynomial identity in~\eqref{eq:budget-evaluation} holds exactly when
\begin{equation}
 \sum_{i<N}\left(e_{c_i}-\sum_{p\in P_i}e_p\right)=e_1
 \quad\text{in the free abelian group on }\mathbb N.
 \label{eq:budget-flow}
\end{equation}
\end{budgetlemma}
\begin{proof}
The distinct monic linear factors $X+j$, $j\in\mathbb N$, have distinct
roots.  Their multiplicities in an identity of products are equal.
This can be proved without invoking a general factorization theorem:
cancel the common copies of $X+j$ from the two sides and evaluate at
$X=-j$; a remaining copy on only one side would make exactly that side
zero.  This applies also to $j=0$, whose factor is $X$.
The resulting multiplicity equations are precisely
\eqref{eq:budget-flow}.
\end{proof}

\subsection{Why cycles cannot manufacture a refutation}

Let $\mathcal R$ be any collection of single-conclusion rules with at most
two premises and labels in $\mathbb N$.  Let $D$ be its least set of
derivable labels, starting from the zero-premise rules.

\begin{budgetlemma}[Single-conclusion flow]
Label $1$ is derivable if and only if some finite list of instances from
$\mathcal R$, with repetitions permitted, satisfies
\eqref{eq:budget-flow}.
\end{budgetlemma}
\begin{proof}
Unroll a finite proof into a tree.  Every non-root node contributes once
as a conclusion and once as a premise of its parent, leaving only the
root vector $e_1$.

Conversely, for a proposed finite list, form the finite closure of just
the rules in that list.  If $1$ is outside this closure, sum the coordinates
of the alleged flow equation that correspond to labels outside the closure.
Whenever a conclusion is outside, at least one premise is outside, since
the set is closed under the listed rule.  The contribution of that rule
to this sum is therefore at most zero.  When its conclusion is inside,
its contribution is again at most zero.  The total on the left is at most
zero, whereas the right side contributes one.  This is impossible.
The finite closure also gives a constructive proof DAG, which can then be
unrolled to a finite tree.
\end{proof}

Repeated premises, self-loops, and disconnected cyclic components do not
invalidate the proof.  The single-conclusion condition is essential: a
rule $a\longrightarrow(a,1)$ would have net vector $e_1$ without a way to
derive either conclusion from no axioms.

\begin{budgettheorem}[Two-accumulator certificate]
Label $1$ is derivable in $\mathcal R$ if and only if there exist $N,M$,
a list of $N$ valid rule instances $P_i\longrightarrow c_i$, and natural
sequences $S_0,\ldots,S_N$, $T_0,\ldots,T_N$ such that
\begin{align}
 S_0&=M+1, & S_N&=1, & T_0&=1, & M&=T_N+1,                \label{eq:budget-ends}\\
 S_{i+1}(M+c_i)&=S_i\prod_{p\in P_i}(M+p),
 &&&& (0\leq i<N),\nonumber\\
 T_{i+1}&=T_i\left(\prod_{p\in P_i}(1+p)+c_i\right),
 &&&& (0\leq i<N).                                       \label{eq:budget-transitions}
\end{align}
No external bound on the labels is required.
\end{budgettheorem}
\begin{proof}
For soundness, the budget recurrence implies $T_i\geq1$ and hence $M\geq2$.
Every factor $M+j$ is positive.  Starting from $S_0=M+1>0$, the first
recurrence shows inductively that every $S_i$ is positive.  Multiply all
of these recurrences and cancel the nonzero internal accumulators.
The endpoints give
$C(M)=(M+1)Q(M)$.
The coefficient-budget lemma, linear-factor uniqueness, and the flow
lemma now prove derivability of $1$.

For completeness, start with a finite proof tree.  List its nodes in a
top-down order in which the next expanded node is on the current frontier.
First compute the entire sequence $T_i$ from this list of labels; set
$M=T_N+1$ only afterwards.  Let $S_i$ be the product of $M+j$ over the
open frontier at stage $i$.  Initially the frontier is just label $1$;
at the end it is empty.  Expanding a node labelled $c_i$ removes its
factor $M+c_i$ and inserts the factors for its actual premises.  This
gives both recurrences and all four endpoints.  A zero-labelled leaf is
expanded with its actual zero-premise truth axiom when that axiom is used.
There is no circular choice of a scale or upper bound.
\end{proof}

The empty list cannot certify target $1$: it would require simultaneously
$M=2$ and $S_0=S_N=1=M+1$.  The argument is entirely finitary and uses
ordinary induction, integer arithmetic, and a finite closure construction.
It can be carried out within ZFC, without a standard-model soundness
assumption or an assumption that ZFC is consistent.

\subsection{The local arithmetic rows are quadratic after branch sharing}

Write the current and next accumulator values as $S,S',T,T'$, and use
$p,q,c$ for the two premise slots and conclusion.  The local flow and
budget constraints are as follows:
\begin{center}
\begin{tabular}{lll}
\hline
Rule arity & Flow equation & Budget equation\\
\hline
Zero & $S'(M+c)=S$ & $T'=T(1+c)$\\
One & $S'(M+c)=S(M+p)$ & $T'=T(1+p+c)$\\
Two & $S'(M+c)=S(M+p)(M+q)$
    & $T'=T(1+p+q+pq+c)$\\
\hline
\end{tabular}
\end{center}
For the explicit truth axiom $c=0$, the first row specializes further to
$S'M=S$ and $T'=T$.

Only the binary row needs additional local arithmetic wires.  Introduce
\[
 u=S(M+p),\qquad w=pq.
\]
Then its four residuals are
\begin{equation}
 u-S(M+p),\quad S'(M+c)-u(M+q),\quad w-pq,\quad
 T'-T(1+p+q+w+c).                                        \label{eq:budget-four-rows}
\end{equation}
They are all quadratic.  The witnesses $u,w$ are uniquely determined and
natural whenever the displayed input coordinates are natural.

The supplied local JSON augments the literal branch rows of
source~68's \texttt{horn\_shared\_source.json}.  Its binary modus-ponens
branch originally uses no scratch wires, so $u,w$ occupy its existing
shared slots \texttt{q0}, \texttt{q1}.  No active branch loses a wire or
shares two distinct values in one slot.  The most demanding inherited
branch still uses thirteen slots.  The generalization branch retains its
own original \texttt{q0}; it needs no new arithmetic wire.  The ten
zero-premise branches, including truth, need no new arithmetic wire either.

For each branch, add the displayed residual squares to that branch's
existing sum of squares.  Multiply the twelve resulting nonnegative
branch polynomials.  Over the naturals it vanishes exactly when one valid
branch and its corresponding accumulator update hold.  The construction
has twenty local existential coordinates:
\[
 p,q,c,\quad a,b,c_{\rm schema},d,\quad q_0,\ldots,q_{12},
\]
where the JSON uses \texttt{v} for the conclusion to avoid the displayed
name collision.  Its five exposed coordinates are $S,S',T,T',M$.
The truth factor rises from degree two to degree four, the inherited
bundle factor stays at degree eight, and all other factors have degree
four.  The local polynomial therefore has exact degree
$11\cdot4+8=52$.  This follows both from exact sparse expansion of every
factor and from nonvanishing of products over $\mathbb Z$.

The following comparison uses exactly the same inherited Horn rows and
the same unflattened branch-product representation.  It does not compare
the size of two final Diophantine polynomials.
\begin{center}
\begin{tabular}{lrr}
\hline
Source-interface quantity & Source 68 & This construction\\
\hline
Local existential coordinates & 23 & 20\\
Exposed local coordinates & 4 & 5\\
Local polynomial degree & 50 & 52\\
Outer certificate coordinates & 5 & 4\\
Bounded universal families & 2 & 1\\
Beta values needed at one step & 2 & 4\\
\hline
\end{tabular}
\end{center}
The extra beta values are a real cost: they encode the second accumulator.
Their arithmetic graph witnesses must be counted before asserting a net
improvement in a compiled polynomial.

\subsection{A completely specified trace interface}

Use four outer natural coordinates $N,x,y,M$ and
\[
 \beta(x,y,j)=x\bmod(1+(j+1)y),\qquad
 S_i=\beta(x,y,2i),\qquad T_i=\beta(x,y,2i+1).
\]
Impose the four endpoint equations~\eqref{eq:budget-ends}.  For every
$i<N$, require a zero of the supplied local polynomial with exposed
coordinates $(S_i,S_{i+1},T_i,T_{i+1},M)$.  This is one bounded universal
family; there is no additional family expressing divisibility of a scale
by all labels.

Every finite pair of accumulator sequences has such a common beta code.
For example, choose $y$ larger than all sequence entries and divisible by
$1,2,\ldots,2N+1$.  The moduli $1+(j+1)y$, $0\leq j\leq2N+1$, are
pairwise coprime: a common divisor is coprime to $y$, divides the index
difference, and that difference divides $y$.  The Chinese remainder
theorem supplies $x$, and the choice of $y$ makes the desired residues
their least nonnegative remainders.  This choice is made after $M$ and
the two traces; it creates no circular scale dependency.

For the calibrated C222 rule relation, the resulting finite-certificate
predicate is equivalent to existence of a C222 refutation of label $1$.
Composing with that calculus's inherited syntactic calibration gives the
same consistency target.  This route checks rule instances directly; its
source does not simulate K215 and does not use a Jones finalizer.

The remaining work is precise.  Substitute literal remainder graphs for
the four beta values, apply a certified bounded-universal elimination to
the unflattened local source, quantify the four outer coordinates, and
only then choose a final degree-reduction or variable-compression
strategy.  The full ordinary polynomial, its complete variable count,
and its coefficient payload have not been produced here.  The theorem
above is not a substitute for those artifacts.

\subsection{An optional normalized variant}

If one wishes to retain the earlier free-zero convention, use weights
$1+jX$ and pad every rule to two premises.  Keep
$T'=T((1+p)(1+q)+c)$ and $M=T_N+1$.  The same coefficient proof applies
to products of $1+jX$, and factor uniqueness recovers all positive-label
multiplicities; the factor for zero is one.  The soundness theorem then
uses zero as a designated derivable label.

The flow equation is
\[
 S'(1+Mc)=S(1+Mp)(1+Mq).
\]
It has a useful exact quadratic lift with four natural auxiliary values:
\[
 w=pq,\quad u=MS,\quad u'=MS',\quad v=Mw.
\]
The six quadratic residuals are
\begin{gather*}
 w-pq,\quad u-MS,\quad u'-MS',\quad v-Mw,\\
 S'+cu'-S-u(p+q+v),\qquad T'-T(1+p+q+w+c).
\end{gather*}
Their sum of squares is quartic, and its zeros are exactly the two desired
updates together with the uniquely determined natural auxiliaries.
Expanding the fifth residual gives
$S'(1+Mc)-S(1+Mp)(1+Mq)$ exactly.  This optional construction explains
how to reuse the previous padded Horn-flow interface, but the unnormalized
linear-factor source is cheaper locally because it respects rule arity.

\subsection{Checks and a necessary negative control}

The standalone checker reconstructs all twelve branch sum-of-squares
polynomials exactly from their literal residuals.  It checks 178 exact
natural branch fixtures against independently written C222 constructor
formulas, including both signs of the guarded bundle.  It exhausts all
60,880 lists of up to three zero-, one-, or two-premise rules on labels
$\{0,1,2\}$, checking the coefficient bound, numerical identity,
polynomial identity, flow equation, and derived-target implication.
Four toy proof-tree fixtures cover repeated premises and explicit truth
leaves; they are finite artificial Horn systems, not ZFC refutations.
The optional normalized lift has its own exact arithmetic checks.

The budget endpoint is essential.  Take $M=11$ and the two rules
$10\longrightarrow3$ and the axiom $7$.
The target $1$ is underivable, but
\[
 (M+3)(M+7)=(M+1)(M+10),\qquad 14\cdot18=12\cdot21.
\]
Indeed $S=(12,18,1)$ satisfies both local flow updates, and $M$ is larger
than every label.  The budget is $T=(1,14,112)$, so the required endpoint
would be $M=113$, not $11$.  The complete checker rejects the fake.
Thus an unspecified assertion that the evaluation point is ``large''
cannot replace the coefficient proof.

These finite checks are diagnostic evidence for the implementation.
The general mathematical justification is the preceding proof, and no
kernel-checked formalization of this new source is claimed.


\section{Coefficient specification and reproducibility}
\label{sec:reproduce}

The display letters for fixed coefficients are abbreviations for actual
integers, not additional quantified unknowns and not a choice depending on
the truth of consistency. The archive contains their complete compressed
little-endian byte representations, exact sizes and SHA-256 digests. The
integer component also reconstructs its coefficient tuple from the supplied
943-tile list and fixed machine table and compares every byte.

The integer equations use the whole
\texttt{mixed-parity-unrounded-1064921} tuple:
$\beta=2^{1064921}-1$, the inherited $K_0,M_C,M_F$, and
the determinant scales $2^{1113}$ and $2^{666104}$. Their six minimal
unsigned payloads total 615,883 raw bytes. The positive equations use the
different whole compact44 tuple, with $\beta=2^{590069}-1$.
The two tuples are never intermixed. In each, $M_F$ is already the shifted
source mask; applying a second shift would specify a different equation.

\subsection{Files that define the equations}

The integer packets under \texttt{integer/generated/} give the ordered
unknowns, residuals, fully expanded abbreviations and source restoration
maps. For a packet with residuals $r_1,\ldots,r_s$, its single equation
is exactly $\sum_{i=1}^s r_i^2=0$. Sparse expanded polynomial files are
provided as additional representations over the fixed coefficient ports.
Substituting the supplied fixed integers gives the required element of
$\Z[x_1,\ldots,x_n]$.

The positive packets are finite straight-line programs in the three binary
operations $+,-,\times$. Their register names introduce no additional
existential variables. Every operation is listed, and a designated final
register defines the polynomial. There is no variable exponent, implicit
division, unbounded product or uninterpreted predicate in any finished
equation.

\subsection{Checks and their mathematical scope}

Run \texttt{python3 verify.py} at the archive root. The normal verification
path uses the Python standard library. It checks authentic payloads,
regeneration of all equation packets, exact polynomial identities,
residual restoration, variable liveness, and the degree witnesses stated
in the proofs. Optimized Python runs of the components were also checked;
the checks do not disappear when \texttt{assert} is disabled.

The additional \texttt{independent\_integer\_audit.py} uses SymPy and
does not import the construction program or its polynomial engine. It
independently restores every row of the joint-guard source and checks
the exact lower-bound coefficients. An independent modular univariate
specialization of the positive circuits supplies another check on their
literal degree. These checks address transcription and algebra. The
universal existence arguments are the written proofs, and the inherited
native-source semantic theorems remain explicitly cited inputs.

The direct component includes positive proof-tree examples, repeated and
zero-premise cases, cycle controls, and a false large-base product collision.
They test the source implementation. The general soundness argument uses
coefficient bounds and finite Horn closure, not successful examples.

For the document, run \texttt{pdflatex report.tex} twice or use the supplied
\texttt{Makefile}. The archive carries source pins, notices and licenses
for all reused material. It includes a content manifest; intermediate
LaTeX files and Python caches are excluded from the release archive.

\section{Further research questions}
\label{sec:questions}

\begin{enumerate}
\item \textbf{Can the coefficient-budget source be compiled efficiently?}
The next decisive deliverable is a sharing-preserving arithmetic
representation of its single bounded family of local checks, followed by
a measured fixed-coefficient polynomial. The local witness count is not
the final witness count. The generic bounded-universal compiler may
reintroduce factorial or power representations internally; the present
source theorem removes its own extra common-multiple block only.

\item \textbf{Can 24 unrestricted variables be reached at degree 8?}
The remaining degree-5 residuals in the decic have explicit origins:
the centered main norm, the strong ordinate, and the centered auxiliary
norm. A further integral coordinate change must control all three at once.
Reducing a formal degree before a required substitution is not enough.

\item \textbf{Can the 23-variable companion have degree below 16?}
Its degree comes from the square of the new dyadic term in the main
residual. Another order of elimination or a different coordinate used to
carry the dyadic divisibility could reduce it. Any restriction must have
a proved lift from authentic canonical source witnesses.

\item \textbf{Can four Pell radicals replace another carrier variable?}
Appendix~\ref{app:four} proves an unconditional separation fact for four
radicals in the positive construction. The open difficulty is preserving
two carrier congruences with one integer coordinate. Their denominators
share powers of $q$, so a naive rational-sum merger allows cancellation
and does not prove the desired divisibilities.

\item \textbf{Can the native congruence move to a different root?}
Supplying $F,\alpha,C$ directly makes positivity cheap, but the constraint
$q\equiv\beta+1\pmod{3\beta}$ still needs to be enforced. Inserting it
into the first root's rational offset can corrupt the transport congruence
when denominators share factors. A useful solution would preserve the
integer strong-root coordinate and allow the two-root descent to apply.

\item \textbf{Can the whole proof become internally checked?}
The latest repository already contains a closed numeric sealed-ZFC
refutation checker. Its internal interpretation and the full native
compiler calibration are separate formalization projects. The finite
algebraic lemmas here are suitable small targets; they should not be
confused with the much larger final consistency biconditional.
\end{enumerate}

\appendix
\section{A four-radical separation lemma for a possible next merger}
\label{app:four}

This appendix records a proved ingredient for the open direction above.
It does not define a new ten-positive-variable consistency equation.
Use the eleven-positive construction's registers, with $I=ic^2$, and set
\[
 t=\sqrt{\ell(\ell+k)+1},\quad
 g=\sqrt{1+\Delta c^2},\quad
 f=\sqrt{1+\Delta I^2},\quad
 y=\sqrt{\frac{QV^2-1}{Q-1}}.
\]
These are real positive roots for every positive input, before imposing
any polynomial equation. Let $b=cT-R$, so $V=f^2b-c$.

\begin{lemma}
Every nonempty signed sum of a subset of $t,g,f,y$ is nonzero.
Consequently, if a signed sum of all four is rational, then all four
roots are rational.
\end{lemma}
\begin{proof}
Here $a\ge2$, $c\ge3$, $I\ge c^2$ and
\[
 t<\ell+k<(a+1)c<g<(a+2)c,
 \qquad f>(a+1)c^2>g+t.
\]
As in the three-root construction, $|V|>1$ and
$y^2=V^2+(V^2-1)/(Q-1)$. The three integer radicands for $t,g,f$
are strictly increasing and less than $Q$. If $y^2$ equalled any one
of them, say $d$, then $Q(d-V^2)=d-1$ with $0<d-1<Q$, impossible.
Thus all four roots are pairwise distinct.

If $b\ne0$, then $y>|V|\ge f^2-c>2f>f+g+t$.
Together with $f>g+t$ and $g>t$, this prevents every signed cancellation.
If $b=0$, then $V=-c$ and $Q>c^2$, so $y<c+1$. Moreover
\[
 (a+1)c-(\ell+k)-(c+1)
   =k(Y^2-1)+a\eta-1>0,
\]
which gives $g>t+y$. Finally
\[
 g+t+y<(2a+4)c+1<(a+1)c^2<f.
\]
The last strict inequality in the middle follows from $a\ge2,c\ge3$.
Thus a subset containing $f$ cannot cancel; without $f$, a subset
containing $g$ cannot cancel; and the remaining two distinct positive
roots $t,y$ cannot cancel.

For rationality, work in the finite multiquadratic field generated by
the four roots. A rational signed sum is fixed by every automorphism.
Subtract its image under any automorphism. If that automorphism changes
any root's sign, the difference is twice a nonempty signed subset sum,
which has just been excluded. Every automorphism therefore fixes every
root. The fixed field is $\Q$, proving the claim.
\end{proof}

The lemma resolves radical separation, not carrier compatibility.
The first transport modulus $q^3(q-1)$ and the main carrier modulus
$4q^2Y$ share the factor $q^2$. Rationality of a sum of quotients with
these denominators does not imply their separate integrality. An attempted
ten-variable equation must supply a new argument at exactly this point.

\begin{thebibliography}{20}
\bibitem{frontier} ProveIt, \emph{Wave-ten proposals and the current targets},
\href{https://github.com/VladimirReshetnikov/ProveIt/blob/2605669a532342f153b13eb6b4a0af4b54f5fa37/Computability/DiophantineConsistency/PROPOSAL_REVIEW_WAVE10.md}{pinned repository review}, 2026.
\bibitem{source61} ProveIt, source 61, \emph{Congruence carriers},
\href{https://github.com/VladimirReshetnikov/ProveIt/tree/2605669a532342f153b13eb6b4a0af4b54f5fa37/Computability/DiophantineConsistency/con-zfc-diophantine-forms/61-congruence-carriers}{pinned package}, 2026.
\bibitem{source64} ProveIt, source 64, \emph{Structural reductions},
\href{https://github.com/VladimirReshetnikov/ProveIt/tree/2605669a532342f153b13eb6b4a0af4b54f5fa37/Computability/DiophantineConsistency/con-zfc-diophantine-forms/64-structural-reductions}{pinned package}, 2026.
\bibitem{source68} ProveIt, source 68, \emph{Joint positivity and free Lucas coordinates},
\href{https://github.com/VladimirReshetnikov/ProveIt/tree/2605669a532342f153b13eb6b4a0af4b54f5fa37/Computability/DiophantineConsistency/con-zfc-diophantine-forms/68-joint-positivity}{pinned package}, 2026.
\bibitem{source69} ProveIt, source 69, \emph{Cones and Pell descent},
\href{https://github.com/VladimirReshetnikov/ProveIt/tree/2605669a532342f153b13eb6b4a0af4b54f5fa37/Computability/DiophantineConsistency/con-zfc-diophantine-forms/69-cones-pell}{pinned package}, 2026.
\bibitem{source70} ProveIt, source 70, \emph{Residue orientation and root transfer},
\href{https://github.com/VladimirReshetnikov/ProveIt/tree/2605669a532342f153b13eb6b4a0af4b54f5fa37/Computability/DiophantineConsistency/con-zfc-diophantine-forms/70-residue-orientation}{pinned package}, 2026.
\bibitem{source71} ProveIt, source 71, \emph{Calibrated determinant guards},
\href{https://github.com/VladimirReshetnikov/ProveIt/tree/2605669a532342f153b13eb6b4a0af4b54f5fa37/Computability/DiophantineConsistency/con-zfc-diophantine-forms/71-determinant-guards}{pinned package}, 2026.
\bibitem{objectzfc} ProveIt, \emph{Object-language ZFC boundary and complete numeric refutation checker},
\href{https://github.com/VladimirReshetnikov/ProveIt/blob/2605669a532342f153b13eb6b4a0af4b54f5fa37/Computability/DiophantineConsistency/OBJECT_ZFC.md}{OBJECT\_ZFC.md} and
\href{https://github.com/VladimirReshetnikov/ProveIt/blob/2605669a532342f153b13eb6b4a0af4b54f5fa37/Computability/DiophantineConsistency/ANNOTATED_REFUTATION_PR.md}{ANNOTATED\_REFUTATION\_PR.md}, 2026.
\bibitem{sun2021} Zhi-Wei Sun, \emph{Further results on Hilbert's Tenth Problem},
Science China Mathematics 64 (2021), 281--306,
\href{https://arxiv.org/abs/1704.03504}{arXiv:1704.03504}, version 7.
\bibitem{sun2026} Zhi-Wei Sun, \emph{Ten unknowns for Hilbert's tenth problem over the integers},
\href{https://arxiv.org/abs/2609.01594}{arXiv:2609.01594}, version 1, 2026.
\bibitem{bayer2026} Jonas Bayer, Marco David, Malte Ha\ss ler, Yuri Matiyasevich,
and Dierk Schleicher, \emph{Diophantine equations over $\Z$: universal bounds and parallel formalization},
\href{https://arxiv.org/abs/2506.20909}{arXiv:2506.20909}, version 3, 2026.
\end{thebibliography}
\end{document}
